\scanpage{155}
\mspage{29}
% archive sheet number: 154

Un précouple de Quillen clos est aussi appelé un
couple de Quillen.

\uline{Proposition} 4.2. Soit $(\Phi,\Psi)$ un précouple de Quillen
dans $M$, et soit $j\in\mathrm{Fl}\,M$.$^{*}$ On a alors les conditions
équivalentes sur $j$ :
% slanted
\margin{NB On utilise seulement que $\Phi\longleftrightarrow\Psi$, et que $j$ puisse se mettre sous la forme $pi$, avec $i\in\Phi$, $p\in\Psi$.}

a) $j\in\widetilde\Phi$

b) $j\in\Psi^*$

c) $j$ est facteur direct d'un $i\in\Phi$

c$'$) $j$ est facteur direct d'un $i\in\Phi$ ayant
même source que $j$, plus précisément il
existe un diagramme

\diagram{\begin{tikzcd}[ampersand replacement=\&,column sep=large]
\& A \arrow[r,"{\unsure{\Delta|j}}"] \arrow[d,"\mathrm{id}_A"'] \& B \arrow[d,"q"']\\
\text{(4.2.1.)} \& A \arrow[u,bend right=36,dashed,"\mathrm{id}_A"'] \arrow[r,"i\unsure{{}^{*}}"] \& B' \arrow[u,bend right=30,dashed,"p"']
\end{tikzcd}}
avec $i\in\Phi$, et
\[
qj=i,\quad pi=j,\quad pq=\mathrm{id}_B,
\tag*{(4.2.\unsure{2|1})}
\]
de sorte que $j$ est un facteur direct de $i\in\Phi$,
par $(\mathrm{id}_A,q)$ admettant la rétraction $(\mathrm{id}_A,p)$.

d) Pour tout diagramme commutatif traits pleins
\diagram{\begin{tikzcd}[ampersand replacement=\&,column sep=small]
\& A \arrow[rr,"j"] \arrow[dr,"i"'] \& \& B \arrow[dl,bend left=65,dotted,"q"]\\
\text{(4.2.3.)} \& \& B' \arrow[ur,"p"'] \&
\end{tikzcd}}
i.e. tte factorisation de $j$ en $pi$, si on a
$p\in\Psi$, il existe une section $q$ de $B'$ sur $B$ \inserted{dans $\unsure{{}_{A\backslash}M|A\backslash M}$},
i.e une flèche $q:B\longrightarrow B'$ telle que
\[
pq=\mathrm{id}_B\quad\text{et}\quad qj=i.
\]

Démonstration de la proposition. Elle se fait par la
suite d'implications suivantes, toutes ($\pm$) tautologiques :
\[
c'\Longrightarrow c\Longrightarrow a\Longrightarrow b\Longrightarrow d\Longrightarrow c',
\]
